%%% FIELD proof-observable-quotient
Fix \(\theta\in\mathcal M^{\circ}_{\mathrm{term},p,q}(G,G)\cap
\mathcal M^{\circ}_{\mathrm{branch},p,q}(G,G)\), and first use the representative of
\(\mathcal R\) generated by \(\theta\).  By \texttt{lem:covariance-branches}, for
each \(k\in V\),
\[
 A^{(k)}-A^{(0)}=a_kh_k^{\mathsf T}.
\]
The nonzero rank minors and factor-line discriminants defining the branch-generic
locus make \(D_k\) nonzero.  Indeed, at a source the displayed formula in
\texttt{lem:covariance-branches} reduces to
\(D_k=(t_k-v_k)a_ka_k^{\mathsf T}\), where \(t_k-v_k\ne0\) on that locus, and at
a non-source the same lemma gives two nonproportional rank-one directions.  Since
every term in the covariance-difference formula contains \(a_k\), in either case
\(D_k\ne0\) implies \(a_k\ne0\).  Choose an observed coordinate \(r(k)\) with
\((a_k)_{r(k)}\ne0\).  Then the preceding rank-one identity determines the row
entry by entry:
\[
 h_{ki}=\frac{(A^{(k)}-A^{(0)})_{r(k),i}}{(a_k)_{r(k)}}
 \qquad(i\in V).
\]
Thus the row recovered in \texttt{def:recovered-closure} is the actual row in
\texttt{def:actual-rows}.  It also satisfies the imposed normalization
\(h_{kk}=0\): by topological ordering, \(P_{jk}=0\) for
\(j\in\operatorname{pa}_G(k)\), and hence
\[
 h_{kk}=\sum_{j\in\operatorname{pa}_G(k)}\delta_{kj}P_{jk}=0.
\]

It remains to determine the support of the recovered row without assuming away
path cancellation.  Apply the terminal-linkage conclusion of
\texttt{thm:terminal-certificate} with \(R=\{k\}\) and \(C=\{i\}\).  The
one-by-one matrix \(H[R,C]\) has rank one exactly when there is a directed path
from \(i\) to \(k^{\star}\) in \(G^{\mathrm{term}}\).  By
\texttt{def:terminal-copy}, such a path exists exactly when there is a
positive-length directed path from \(i\) to \(k\) in \(G\): a path to
\(k^{\star}\) ends with \(j\to k^{\star}\) for a parent \(j\) of \(k\), and
replacing that last edge by \(j\to k\) gives the path in \(G\), while the reverse
replacement gives the converse.  Consequently,
\[
 (H_{\mathcal R})_{ki}\ne0
 \quad\Longleftrightarrow\quad
 i\to k\in\overline G.
\]
The definition of \(\mathcal C(\mathcal R)\) therefore yields
\(\mathcal C(\mathcal R)=\overline G\) for this representative.

We next prove representative invariance.  The monomial group is generated by
nonsingular diagonal matrices and permutation matrices.  Let
\(C=\operatorname{diag}(c_1,\ldots,c_q)\), with every \(c_i\ne0\).  Under the
diagonal action in \texttt{def:aligned-factor-input},
\[
 A^{(k)}C-A^{(0)}C
 =a_kh_k^{\mathsf T}C
 =(c_ka_k)\bigl(c_k^{-1}h_k^{\mathsf T}C\bigr).
\]
Uniqueness of the row just proved shows that the recovered matrices before and
after rescaling satisfy
\[
 \widetilde H=C^{-1}HC.
\]
It follows entrywise that \(\widetilde H_{ki}=c_k^{-1}H_{ki}c_i\), so the zero
pattern is unchanged.  If \(U'=I_q-\Lambda'^{(0)}\) is a candidate structural
witness, put \(\widetilde U=C^{-1}U'C\).  Then
\[
 (\widetilde H\widetilde U)_{ki}
 =c_k^{-1}(HU')_{ki}c_i.
\]
Thus every exact-support and noncancellation condition is preserved.  Likewise,
for all index sets \(R_0,C_0\subseteq V\),
\[
 \operatorname{rank}\widetilde H[R_0,C_0]
 =\operatorname{rank}\bigl(C[R_0,R_0]^{-1}H[R_0,C_0]C[C_0,C_0]\bigr)
 =\operatorname{rank}H[R_0,C_0],
\]
because both diagonal submatrices are invertible.  These equalities transport the
terminal-rank, affine-feasibility, and exact-support certificates of
\texttt{thm:terminal-certificate} in both directions.

For completeness, consider any model witness
\((F',\Lambda'^{(0)},\delta',v',t')\in\mathcal M_{p,q}(D)\) for a candidate DAG
\(D\).  Define, in every context,
\[
 \widetilde F=F'C,
 \qquad
 \widetilde\Lambda^{(k)}=C^{-1}\Lambda'^{(k)}C,
 \qquad
 \widetilde\varepsilon^{(k)}=C^{-1}\varepsilon'^{(k)}.
\]
Since \(C^{-1}e_k=c_k^{-1}e_k\), the intervention remains a single-node
intervention and its row, baseline variances, and target variances are
\[
 \widetilde\delta_k^{\mathsf T}=c_k^{-1}\delta_k'{}^{\mathsf T}C,
 \qquad
 \widetilde v_i=c_i^{-2}v'_i,
 \qquad
 \widetilde t_i=c_i^{-2}t'_i.
\]
Hence all variances remain strictly positive, and
\[
 \widetilde\Lambda^{(0)}_{ki}+\widetilde\delta_{ki}
 =c_k^{-1}\bigl(\Lambda'^{(0)}_{ki}+\delta'_{ki}\bigr)c_i.
\]
The exact edge support and retained-edge inequalities are therefore preserved.
Right multiplication of \(F'\) by an invertible diagonal matrix preserves its
rank, its nonzero columns, and its pairwise noncollinear column lines.  Nonzero
coordinatewise scaling also preserves centering, independence, and
non-Gaussianity of the error coordinates.  Finally,
\[
 (I_q-\widetilde\Lambda^{(k)})^{-1}
 =C^{-1}(I_q-\Lambda'^{(k)})^{-1}C,
\]
and therefore
\[
 \widetilde F(I_q-\widetilde\Lambda^{(k)})^{-1}
 \widetilde\varepsilon^{(k)}
 =F'(I_q-\Lambda'^{(k)})^{-1}\varepsilon'^{(k)}.
\]
In particular the full observed laws, and hence every observed covariance
\(\Sigma^{(k)}\), are unchanged.  The construction with \(C^{-1}\) is its
inverse, so this is a bijection of realization witnesses and of positive-real
covariance-feasibility fibers.

Now let \(\Pi\) be a permutation matrix and define \(\sigma\) by
\(\Pi e_j=e_{\sigma(j)}\).  Keep the baseline label \(0\) fixed and relabel each
target context \(j\in V\) as the old context \(\sigma(j)\).  Set
\[
 \widetilde\Lambda^{(0)}=\Pi^{\mathsf T}\Lambda'^{(0)}\Pi,
 \qquad
 \widetilde\Lambda^{(j)}=\Pi^{\mathsf T}\Lambda'^{(\sigma(j))}\Pi,
 \qquad
 \widetilde F=F'\Pi,
\]
and transform the corresponding errors to \(\Pi^{\mathsf T}\varepsilon'^{(0)}\)
and \(\Pi^{\mathsf T}\varepsilon'^{(\sigma(j))}\), respectively.  Direct
substitution gives the same observed law in the relabeled context and gives
\[
 \widetilde H=\Pi^{\mathsf T}H\Pi.
\]
Thus a common permutation merely relabels the graph, targets, supports, terminal
paths, and all feasibility systems.  Combining this permutation transport with
the diagonal transport proves invariance under every element of
\(\mathfrak M_q\).  A realization witness is carried to the displayed witness;
an infeasibility or separator certificate is carried back by the inverse
bijective substitution.  Consequently feasibility and nonmembership, not merely
their truth values, have representative-independent certificates.

The projective gauge in \texttt{def:aligned-factor-input} removes exactly the
diagonal part.  Indeed, for a nonzero loading column \(b_j\), let \(s(j)\) be its
least-index nonzero observed coordinate and take
\(c_j=b_{s(j),j}^{-1}\).  The \(s(j)\)-th coordinate of \(c_jb_j\) is then one.
If another nonzero scalar gives the same normalization, comparison of that pivot
coordinate shows that it equals \(c_j\).  Thus the scale is unique on each
rational projective chart.  Only the finite common permutation orbit remains,
and the preceding conjugation calculation proves that enumerating it preserves
terminal ranks, supports, covariance feasibility, witnesses, and separators.

By \texttt{def:observable-classes}, membership of a DAG in
\(\mathcal G_{\mathrm{soft}}(\mathcal R)\) or in
\(\mathcal G_{\mathrm{cov}}(\mathcal R,\{\Sigma^{(k)}\}_{k\in K})\) is an
existential assertion about precisely the witnesses transported above.  The
diagonal action leaves each labeled graph unchanged, and the permutation action
relabels every graph simultaneously.  Hence both graph-set maps, and therefore
their set difference, are independent of the representative and define the
claimed quotient-valued pointwise population functional.

Finally, this pointwise conclusion must not be identified with the generic graph
classes.  By \texttt{def:observable-classes}, observable covariance membership
asks for a witness at the single supplied covariance tuple.  In contrast, by
\texttt{def:covariance-class}, \(G'\in\mathsf{Cov}_p(G)\) is a statement about
the closure of \(\mathcal E_p(G,G')\) in every full-dimensional sign and variance
chamber.  The quantifiers are inequivalent: if \(C_0\) is a positive-dimensional
semialgebraic chamber and \(u\in C_0\), then
\(C_0\setminus\{u\}\) is semialgebraic and dense in \(C_0\) but does not contain
\(u\), whereas \(\{u\}\) contains the supplied point but is not dense in
\(C_0\).  These are explicit witnesses that chamberwise genericity alone implies
neither pointwise membership at every point nor its converse.  Similarly,
\texttt{def:corrected-class} quantifies actual-row compatibility over every point
of its stated generic locus, while \(\mathcal G_{\mathrm{soft}}(\mathcal R)\)
asks only for a realization of the supplied \(\mathcal R\).  No assumption or
established result in the dependency list asserts constancy of either feasibility
property across those loci.  Therefore the proved quotient-valued pointwise
difference is distinct from, and is not identified here with,
\(\mathsf{Cov}_p(G)\setminus\mathsf H_{\mathrm{soft}}(G)\), as claimed.
%%% FIELD statement-observable-quotient
On \(\mathcal M^{\circ}_{\mathrm{term},p,q}(G,G)\cap\mathcal M^{\circ}_{\mathrm{branch},p,q}(G,G)\), the relation recovered from the aligned factorization datum satisfies \(\mathcal C(\mathcal R)=\overline G\) up to the common latent permutation.  The maps \(\mathcal R\mapsto\mathcal G_{\mathrm{soft}}(\mathcal R)\) and \((\mathcal R,\{\Sigma^{(k)}\}_{k\in K})\mapsto\mathcal G_{\mathrm{cov}}(\mathcal R,\{\Sigma^{(k)}\}_{k\in K})\) are invariant to the representative of \(\mathcal R\) under \(\mathfrak M_q\).  The rational projective gauge removes every nonzero diagonal scale, and enumeration of the remaining common permutation orbit preserves terminal ranks, exact-support conditions, covariance feasibility, realization witnesses, and separators.  Consequently, \(\mathcal G_{\mathrm{cov}}(\mathcal R,\{\Sigma^{(k)}\}_{k\in K})\setminus\mathcal G_{\mathrm{soft}}(\mathcal R)\) is a well-defined quotient-valued pointwise population functional.  Without an additional generic-to-pointwise constancy theorem, this proposition does not identify that pointwise difference with the chamberwise graph functional \(\mathsf{Cov}_p(G)\setminus\mathsf H_{\mathrm{soft}}(G)\).
